% ============================================================
\section{The test set in full}
\label{sec:testset}

The test set contains $180$ items in five families
(Table~\ref{tab:families}); scoring counts occurrences of a target unit
against an expected count, not string matches against a reference
transcript.

\begin{table*}[p]
\centering
\caption{Item families.}
\label{tab:families}
\begin{tabular}{lrl}
\toprule
Family & Count & What it tests \\
\midrule
word repetition & 60 & a single word repeated $k$ times in a fixed carrier \\
matched control & 54 & the same carrier, $k$ \emph{distinct} filler words (no repetition) \\
sentence repetition & 32 & a whole short sentence repeated $k$ times \\
tongue twisters & 24 & a twister sentence (with internal repetition) repeated $k$ times \\
numbers & 10 & English number phrases with intrinsic digit-word repetition \\
\midrule
Total & 180 & \\
\bottomrule
\end{tabular}
\end{table*}

\paragraph{The $k$-ladder.} Word repetition and controls run
$k \in \{1,2,3,4,$ $6,8,12,16,$ $24,32\}$, bracketing the expected collapse
point $k^\ast \in 4$--$16$; sentence repetition stops at $k=16$, twisters
at $k=4$ (Table~\ref{tab:twisters}); numbers are fixed strings, no
ladder (Table~\ref{tab:numbers}).

\paragraph{Word-repetition carriers.} Table~\ref{tab:templates} reproduces
all six carrier templates verbatim. The six targets are common, short,
phonetically distinct from every filler, and homophone-free, limiting
judge mis-transcription noise; the judge audit (Section~\ref{sec:judge})
later located the transcription risk in the judge's autoregressive
decoder, not the word list.
\begin{table*}[p]
\centering
\caption{Word-repetition carriers: fixed prefix, target repeated $k$ times,
fixed suffix; the filler pool builds the matched control at the same $k$.}
\label{tab:templates}
\small
\begin{tabular}{@{}lp{3.2cm}p{1.4cm}p{4.4cm}p{6.6cm}@{}}
\toprule
ID & Prefix & Target & Suffix & Filler pool (for the control) \\
\midrule
t1 & The dog was & very & big and it ran across the field. & quite, really, truly, fairly, rather, somewhat, extremely, notably \\
t2 & She said & no & to the offer and then left the room. & yes, maybe, sure, fine, okay, well, right, hmm \\
t3 & He walked & far & into the forest before turning back. & deep, fast, long, wide, high, low, near, past \\
t4 & The light was & blue & before the storm arrived. & green, bright, pale, dim, warm, cold, sharp, soft \\
t5 & The runner moved & quick & along the narrow path. & swift, smooth, steady, light, sharp, loose, tight, clean \\
t6 & They were & really & tired after the long journey. & truly, quite, very, rather, deeply, clearly, plainly, surely \\
\bottomrule
\end{tabular}
\end{table*}

\paragraph{Number phrases.} Table~\ref{tab:numbers} reproduces all ten
phrases verbatim. One digit-word recurs at several magnitudes, so correct
rendering needs a place-value counter, not adjacent-token runs.
\begin{table*}[p]
\centering
\caption{Number-phrase items, with the target digit-word and its true count.}
\label{tab:numbers}
\small
\begin{tabular}{@{}lp{12.3cm}ll@{}}
\toprule
ID & Phrase & Target & Count \\
\midrule
n01 & six hundred sixty six thousand six hundred sixty six & six & 6 \\
n02 & seven hundred seventy seven thousand seven hundred seventy seven & seven & 6 \\
n03 & nine hundred ninety nine thousand nine hundred ninety nine & nine & 6 \\
n04 & three hundred thirty three thousand three hundred thirty three & three & 6 \\
n05 & six hundred sixty six million six hundred sixty six thousand six hundred sixty six & six & 9 \\
n06 & eight hundred eighty eight thousand eight hundred eighty eight & eight & 6 \\
n07 & five hundred fifty five & five & 3 \\
n08 & four hundred forty four & four & 3 \\
n09 & two hundred twenty two million two hundred twenty two thousand two hundred twenty two & two & 9 \\
n10 & one hundred eleven thousand one hundred eleven & one & 5 \\
\bottomrule
\end{tabular}
\end{table*}

\paragraph{Sentence repetition and tongue twisters.} Sentence items
(Table~\ref{tab:sentences}) score one content word at expected count $k$;
twister items expect per-copy occurrences $\times\,k$, the w6 and w7
targets (per-copy $0$) probing phonological stress only, outside the
counting metric.
\begin{table*}[p]
\centering
\caption{Sentence-repetition items ($k \in \{1,2,3,4,6,8,12,16\}$,
expected count $k$) and tongue-twister items ($k\in\{1,2,4\}$, expected
count per-copy occurrences of the target $\times\,k$).}
\label{tab:sentences}\label{tab:twisters}
\small
\begin{tabular}{@{}lp{11.5cm}ll@{}}
\toprule
ID & Sentence & Target & Per-copy \\
\midrule
s1 & The bell rang twice. & bell & --- \\
s2 & He counted the stones. & stones & --- \\
s3 & Rain fell on the roof. & rain & --- \\
s4 & The door stayed open. & door & --- \\
\midrule
w1 & She sells sea shells by the sea shore. & sea & 2 \\
w2 & Peter Piper picked a peck of pickled peppers. & peck & 1 \\
w3 & How much wood would a woodchuck chuck. & chuck & 1 \\
w4 & Red lorry yellow lorry red lorry yellow lorry. & lorry & 4 \\
w5 & The sixth sick sheikh's sixth sheep is sick. & sixth & 2 \\
w6 & Six slick slim sycamore saplings. & s (alliteration only) & 0 \\
w7 & Fresh French fried fly fritters. & fr (alliteration only) & 0 \\
w8 & Truly rural truly rural truly rural. & rural & 3 \\
\bottomrule
\end{tabular}
\end{table*}

\paragraph{Matched controls.} Same prefix, suffix, and word count, the
$k$ target copies replaced by $k$ distinct fillers cycled from the pool,
the target never appearing in its own control. Only periodicity is
removed: an item-control difference cannot reflect sequence length,
generation duration, or text amount, the basis of the main text's
dissociation test.

\paragraph{Repetition-boundary lists.} Word- and sentence-repetition
items list their $k$ target copies in order, controls their $k$ fillers;
the boundary-location step (Section~\ref{sec:instrumentation}) searches
one explicit, model-agnostic list, left-to-right and non-rewinding in
both arms, so fitted decay rates compare directly.

\paragraph{The instrumented subset.} Teacher-forced hidden-state and
attention capture (Section~\ref{sec:instrumentation}) covers word
repetition, sentence repetition, and controls at $k\ge2$, $136$ items;
twisters and numbers are behavioural-only.

% ============================================================
\section{Scoring}
\label{sec:scoring}

\noindent\textbf{Two flags.} As built, Count A is Whisper
large-v3~\cite{radford2023whisper}; the judge audit
(Section~\ref{sec:judge}) showed its autoregressive decoder
systematically undercounts correct periodic speech, so the main text uses
a CTC judge scored by relative count error, and Whisper-derived numbers
are flagged where they recur. A judge swap changes only the transcript
source: Count B is recalibrated per judge, and every
transcript-independent piece survives.

\paragraph{Three metrics.} \textbf{Count A} counts exact whole-token
target matches in the normalised transcript (lower-cased, punctuation
stripped, numerals expanded to number words). \textbf{Count B} inverts a
per-model, per-family, per-template linear fit of duration on $k$,
calibrated only on a trusted low-$k$ regime ($k\le3$, Count A exact,
duration $>0.2$\,s) fixed before seeing high-$k$ behaviour. An item is
\textbf{correct} only if Count A exactly equals the expected count
\emph{and} the duration ratio against the fit lies in $[0.70, 1.45]$.

\paragraph{Exclusions.} Audio that is empty or degenerate (duration
$<0.25$\,s, RMS $<10^{-3}$, spectral flatness $>0.35$, or an empty
transcript) carries no count and is reported as its own rate, checked
before any transcript-based label. At
$|\text{Count A} - \text{Count B}| > \max(1,\,
0.15\cdot\text{expected count})$ the item is a flagged disagreement,
never a synthesised number.
